\appendix
\section{Mathematical formulation and verification}
\label{app:mathematical_formulation}

\subsection{Reference mapping and equilibrium}
Let $\Omega_0$ be the FE template and $\boldsymbol\chi_i(\mathbf X)=\mathbf X+\boldsymbol\phi_i(\mathbf X)$ the unloaded subject geometry. The anatomical mapping and the elastic displacement are different fields. With $\mathbf F_i=\mathbf I+\nabla_X\boldsymbol\phi_i$ and $J_i=\det\mathbf F_i$, the physical small-strain tensor is
\begin{equation}
\boldsymbol\varepsilon_i(\mathbf u)=\operatorname{sym}(\nabla_X\mathbf u\,\mathbf F_i^{-1}).
\end{equation}
The pullback requires an orientation-preserving, nonsingular map. RBF smoothness alone does not prove this condition. For physical coordinates in mm, the discrete weak problem is
\begin{equation}
\int_{\Omega_i}\mathbf C_i\boldsymbol\varepsilon(\mathbf u):\boldsymbol\varepsilon(\mathbf v)\,dV
+\int_{\Gamma_{S1,i}}\gamma\,\mathbf u\cdot\mathbf v\,dS
+\mathbf v_h^T\mathbf K_{lig,i}\mathbf u_h
=\int_{\Gamma_{load,i}}\mathbf t\cdot\mathbf v\,dS+\mathbf v_h^T\mathbf f_{pre,i}.
\end{equation}
The surface penalty has units N/mm$^3$, because its associated traction is $\gamma\mathbf u$ in N/mm$^2$. It imposes an approximate fixed boundary, not exact $\mathbf u=0$. The ligament stiffness and pretension terms are required for consistency with the assembled solver. For one spring segment of length $L$, direction $\mathbf e$, assigned stiffness $k$ and pretension $T$, its directional stiffness is
\begin{equation}
\mathbf W=k\mathbf e\mathbf e^T+(T/L)(\mathbf I-\mathbf e\mathbf e^T).
\end{equation}
This matrix is assembled with opposite off-diagonal blocks at its two endpoints. It is a linearized spring representation without tension-only switching.

The surface measure obeys $dS_i=J_{surf,i}dS_0$, where
\begin{equation}
J_{surf,i}=J_i\|\mathbf F_i^{-T}\mathbf n_0\|,\qquad
\mathbf t=\frac{\mathbf F^{tot}}{\int_{\Gamma_{load,0}}J_{surf,i}\,dS_0}.
\end{equation}
Thus the applied physical force integral equals $\mathbf F^{tot}$. The penalty and load integrals use the corresponding mapped surface measures. The force vectors are listed in Table~\ref{tab:loads}; their global axes are not assumed identical to the PCA reporting axes below.

\subsection{Unloaded and loaded surfaces}
Let $\mathbf X_j$ be surface template nodes. The two point clouds used for kinematics are
\begin{equation}
\mathbf p^0_{ij}=\mathbf X_j+\boldsymbol\phi_i(\mathbf X_j),\qquad
\mathbf p^1_{i\ell j}=\mathbf p^0_{ij}+\mathbf u_{i\ell}(\mathbf X_j).
\end{equation}
Both use identical topology and point correspondence. Comparing $\mathbf X_j$ with $\mathbf p^1_{i\ell j}$ would incorrectly include subject morphology in the measured motion. Geometry was reconstructed with the original two-stage interpolation: template-to-FE mapping (10 neighbors, smoothing 50) followed by FE-to-surface mapping (10 neighbors, zero smoothing). Saved elastic fields use the latter interpolation. Cached sparse interpolation weights implement the same linear RBF evaluation; numerical comparisons against direct RBF evaluation check this equivalence.

\subsection{Pelvis-aligned reporting frame}
The origin $\mathbf o$ is the centroid of the unloaded surface. The initial ML axis is the dominant eigenvector of the point-cloud covariance, with deterministic sign. The symmetry-plane routine uses this PCA approximation; it does not optimize reflection symmetry. Project the centered points onto its orthogonal plane and take the leading in-plane singular vector. Orthogonalization gives $\hat{\mathbf z}$, followed by
\begin{equation}
\hat{\mathbf y}=\frac{\hat{\mathbf z}\times\hat{\mathbf x}}{\|\hat{\mathbf z}\times\hat{\mathbf x}\|},\qquad
\hat{\mathbf z}=\hat{\mathbf x}\times\hat{\mathbf y},\qquad
\mathbf P=[\hat{\mathbf x},\hat{\mathbf y},\hat{\mathbf z}].
\end{equation}
The ML and AP axes are simultaneously reversed if needed so that ML points from the left to right ilium centroids. This construction is orthonormal and right-handed but does not independently anchor the AP/CC signs to anatomical landmarks. Absolute components are therefore used for comparisons. No clinical nutation sign is inferred from the stored $\alpha_x$ sign.

\subsection{Affine correction, regions and rigid fits}
Fit the sacral point correspondence by ordinary least squares,
\begin{equation}
(\mathbf A_S,\mathbf b_S)=\arg\min_{\mathbf A,\mathbf b}\sum_{j\in S}
\|\mathbf p^1_j-\mathbf A\mathbf p^0_j-\mathbf b\|^2,
\qquad \mathbf q_j=\mathbf A_S^+(\mathbf p^1_j-\mathbf b_S),
\end{equation}
where $+$ denotes the Moore--Penrose pseudoinverse. This removes fitted affine sacral deformation, including stretch and shear; it is not only a change of rigid observer frame.

For each side, select sacral and iliac points within 5~mm of the opposing point cloud. If fewer than 80 points qualify, use up to the nearest 80 available points. Distances are between discrete points rather than continuous articular surfaces. For each bone ROI, the rigid fit maps $\mathbf p^0$ to $\mathbf q$. With centered row matrices $\widetilde{\mathbf P}_B,\widetilde{\mathbf Q}_B$,
\begin{equation}
\widetilde{\mathbf P}_B^T\widetilde{\mathbf Q}_B=\mathbf U\boldsymbol\Sigma\mathbf V^T,
\quad \mathbf R_B=\mathbf V\operatorname{diag}(1,1,\det(\mathbf V\mathbf U^T))\mathbf U^T,
\quad \mathbf t_B=\bar{\mathbf q}_B-\mathbf R_B\bar{\mathbf p}^0_B.
\end{equation}
The fit is repeated after excluding the largest 10\% residuals, using two trimming iterations. Trimming reduces outlier influence but does not validate the ROI as a physical contact patch.

\subsection{Relative rotation and translation}
For sacral and iliac fits $S,I$,
\begin{equation}
\mathbf R_{rel}=\mathbf R_S^T\mathbf R_I,\qquad
\mathbf R_{loc}=\mathbf P^T\mathbf R_{rel}\mathbf P
=\mathbf R_z(\gamma_z)\mathbf R_y(\beta_y)\mathbf R_x(\alpha_x).
\end{equation}
The fixed-axis $xyz$ extraction, for zero-based matrix indices away from gimbal lock, is
\begin{equation}
\alpha_x=\operatorname{atan2}(R_{21},R_{22}),\quad
\beta_y=\arcsin[-\operatorname{clip}(R_{20},-1,1)],\quad
\gamma_z=\operatorname{atan2}(R_{10},R_{00}).
\end{equation}
At $\beta_y=\pm\pi/2$, choose $\gamma_z=0$ and
$\alpha_x=\operatorname{atan2}[-\operatorname{sign}(R_{20})R_{01},R_{11}]$.
In implementation, the equivalent stable expression $\beta_y=\operatorname{atan2}(-R_{20},\sqrt{R_{00}^2+R_{10}^2})$ is used, with the singular branch restricted to vanishing cosine. The singular branch is included for completeness; physiological small-angle outputs should remain far from it. Matrix reconstruction tests cover both branches.

For the unloaded sacral ROI centroid $\mathbf p_c$, the reported displacement is
\begin{equation}
\mathbf d=\mathbf P^T[(\mathbf R_I-\mathbf R_S)\mathbf p_c+\mathbf t_I-\mathbf t_S].
\end{equation}
It evaluates the difference of two fitted maps at a common reference point. It is neither the distance between opposing articular surfaces nor a displacement along a surface normal. The bilateral summaries are
\begin{align}
|\boldsymbol\theta|&=\tfrac12\sum_{s=L,R}(\alpha_{x,s}^2+\beta_{y,s}^2+\gamma_{z,s}^2)^{1/2}, &
|\mathbf d|&=\tfrac12\sum_{s=L,R}\|\mathbf d_s\|,\\
|d_k|&=\tfrac12(|d_{k,L}|+|d_{k,R}|), &
A_d&=\big|\|\mathbf d_L\|-\|\mathbf d_R\|\big|.
\end{align}
Angles are converted to degrees before reporting. The exact principal rotation angle would instead be $\arccos[(\operatorname{tr}\mathbf R_{rel}-1)/2]$ with numerical clipping. The Cardan norm approximates it at small angles and depends on the chosen axes and sequence.

\subsection{Size and conditional scaling}
The archived allometry workflow sums volumes of the selected left ilium, right ilium and sacrum surface bodies, using mapped surface geometry and converting mm$^3$ to cm$^3$. These are whole-bone geometric volumes, not birth-canal volume or a direct integral of cortical tissue. The size factor is $s_i=(V_i/V_{template})^{1/3}$. Therefore a log-volume slope $b$ becomes $3b$ under log scale, with identical fitted responses and tests.

For the sex-interaction log model, $b_M=\beta_v$ and $b_F=\beta_v+\beta_{v:s}$. The female slope variance is
\begin{equation}
\operatorname{Var}(b_F)=\operatorname{Var}(\beta_v)+\operatorname{Var}(\beta_{v:s})+2\operatorname{Cov}(\beta_v,\beta_{v:s}).
\end{equation}
Reported intervals use this full covariance. Exponentiating a fitted log response gives a conditional geometric mean, not an unbiased arithmetic mean. Figure~\ref{fig:allometry} uses this same adjusted model, with the other continuous covariates set to their cohort means.
